% from ega2-1-fr.tex
\oldpage[II]{146}
valuation convenable, on peut supposer que les $c_i$ appartiennent tous à
$A$, et que l'un au moins est inversible. Mais alors
(\hyperref[II.4.2.6-fr]{4.2.6}), le système $(c_0,\ldots,c_n)$ définit aussi
un point de $\mathbf{P}_Y^n$ à valeurs dans $A$, ce qui démontre notre
assertion.
\item[\rm (iii)] Les critères (\hyperref[II.7.2.3-fr]{7.2.3}) et
(\hyperref[II.7.3.8-fr]{7.3.8}) sont surtout commodes quand on considère la
donnée d'un $Y$-préschéma $X$ comme équivalente à la donnée du foncteur
\[
X(Y')=\operatorname{Hom}_Y(Y',X)
\]
en un $Y$-préschéma $Y'$~; ces critères nous permettront par exemple de
prouver que sous certaines conditions les «~schémas de Picard~» sont propres.
\end{enumerate}
\end{remark}

\begin{corollary}[7.3.10]
\label{II.7.3.10-fr}
Soient $Y$ un schéma intègre (resp. un schéma intègre localement noethérien),
$X$ un schéma intègre, $f:X\to Y$ un morphisme dominant.
\begin{enumerate}
\item[\rm (i)] Si $f$ est quasi-compact et universellement fermé, tout anneau
de valuation dont le corps des fractions est le corps $R(X)$ des fonctions
rationnelles sur $X$, et qui domine un anneau local de $Y$, domine aussi un
anneau local de $X$.
\item[\rm (ii)] Inversement, supposons que $f$ soit de type fini, et que la
propriété énoncée dans (i) soit vérifiée par tout anneau de valuation (resp.
tout anneau de valuation discrète) ayant $R(X)$ pour corps des fractions.
Alors $f$ est propre.
\end{enumerate}
\end{corollary}

Notons d'abord que les hypothèses impliquent dans tous les cas que $f$ est
séparé (\hyperref[I.5.5.9-fr]{I, 5.5.9}).
\begin{enumerate}
\item[\rm (i)] Soient $K=R(Y)$, $L=R(X)$, $y$ un point de $Y$, $A$ un anneau
de valuation ayant $L$ pour corps des fractions et dominant
$\mathcal{O}_y$~; l'injection $\mathcal{O}_y\to A$ définit donc un morphisme
$h$ de $Y'=\operatorname{Spec}(A)$ dans $Y$
(\hyperref[I.2.4.4-fr]{I, 2.4.4}) tel que $h(a)=y$, en désignant par $a$ le
point fermé de $Y'$~; en outre, si $\eta$ est le point générique de $Y$, qui
est aussi celui de $\operatorname{Spec}(\mathcal{O}_y)$, on a $h(b)=\eta$, en
désignant par $b$ le point générique de $Y'$ (puisque $K\subset L$ par
hypothèse). Si $\xi$ est le point générique de $X$, on a $k(\xi)=k(b)=L$ par
hypothèse, d'où un $Y$-morphisme $g:\operatorname{Spec}(L)\to X$ tel que
$g(b)=\xi$~; en vertu de (\hyperref[II.7.3.8-fr]{7.3.8}), $g$ provient d'un
$Y$-morphisme $g':Y'\to X$. Si $x=g'(a)$, il est clair que $A$ domine
$\mathcal{O}_x$.
\item[\rm (ii)] La question étant locale sur $Y$, on peut toujours supposer
que $Y$ est affine (resp. affine et noethérien). Comme $f$ est de type fini,
on peut appliquer dans les deux cas le lemme de Chow
(\hyperref[II.5.6.1-fr]{5.6.1}). Il y a donc un morphisme projectif
$p:P\to Y$, un morphisme d'immersion $j:X'\to P$ et un morphisme projectif,
surjectif et birationnel $g:X'\to X$ ($X'$ étant intègre) tels que le diagramme
\[
\begin{tikzcd}[column sep=large,row sep=large]
P \arrow[d,"p"'] & X' \arrow[l,"j"'] \arrow[d,"g"] \\
Y & X \arrow[l,"f"]
\end{tikzcd}
\]
soit commutatif. Il suffit de prouver que $j$ est une immersion \emph{fermée},
car alors $f\circ g=p\circ j$ sera un morphisme projectif, donc propre, et
comme $g$ est surjectif, $f$ sera aussi propre
(\hyperref[II.5.4.3-fr]{5.4.3}). Soit $Z$ le sous-préschéma fermé réduit de
$P$ ayant $\overline{j(X')}$ comme espace sous-jacent
(\hyperref[I.5.2.1-fr]{I, 5.2.1})~; comme $X'$ est intègre, $j$ se factorise
en $i\circ h$, où $i:Z\to P$ est l'injection canonique, $h:X'\to Z$ une
immersion ouverte dominante (\hyperref[I.5.2.3-fr]{I, 5.2.3}), et $Z$ est
intègre~;

